%%%
%%% Radical "⾥"
%%%
\section*{Radical 166: ``⾥''}\addcontentsline{toc}{section}{Radical 166: ⾥}\addcontentsline{loh}{figure}{\#\#\#\# 166: ⾥}

%%%%%%%%%% 里 %%%%%%%%%%
\subsection*{里}\addcontentsline{loh}{figure}{里}

\begin{Entry}{里}{7}{⾥}[Kangxi 166]
  \begin{Phonetics}{里}{li3}[][HSK 1]
    \definition*{s.}{Sobrenome: Li}
    \definition{clas.}{li, uma unidade chinesa de comprimento (= 1/2 quilômetro)}
    \definition{s.}{forro; revestimento; interior; parte de trás do tecido | interno; dentro; no interior | vizinhança; vizinhos | cidade natal; local de origem}
  \seealsoref{上}{shang4}
  \seealsoref{中}{zhong1}
  \synonymref{内}{nei4}
  \synonymref{上}{shang4}
  \synonymref{中}{zhong1}
  \antonymref{边}{bian1}
  \antonymref{表}{biao3}
  \antonymref{面}{mian4}
  \antonymref{外}{wai4}
  \end{Phonetics}
\end{Entry}

\begin{Entry}{里头}{7,5}{⾥,⼤}
  \begin{Phonetics}{里头}{li3tou5}[][HSK 2]
    \definition{s.}{dentro}
  \synonymref{里边}{li3bian5}
  \synonymref{里面}{li3mian5}
  \antonymref{外边}{wai4bian5}
  \antonymref{外部}{wai4bu4}
  \antonymref{外头}{wai4tou5}
  \antonymref{外面}{wai4mian5}
  \end{Phonetics}
\end{Entry}

\begin{Entry}{里边}{7,5}{⾥,⾡}
  \begin{Phonetics}{里边}{li3bian5}[][HSK 1]
    \definition{s.}{em; dentro; no interior}
  \synonymref{里头}{li3tou5}
  \synonymref{里面}{li3mian5}
  \synonymref{内部}{nei4bu4}
  \antonymref{外边}{wai4bian5}
  \antonymref{外头}{wai4tou5}
  \antonymref{外面}{wai4mian5}
  \end{Phonetics}
\end{Entry}

\begin{Entry}{里面}{7,9}{⾥,⾯}
  \begin{Phonetics}{里面}{li3mian5}[][HSK 3]
    \definition{s.}{dentro; interior}
  \end{Phonetics}
\end{Entry}

\begin{Entry}{里斯本}{7,12,5}{⾥,⽄,⽊}
  \begin{Phonetics}{里斯本}{li3si1ben3}
    \definition*{s.}{Lisboa, capital de Portugal}
  \end{Phonetics}
\end{Entry}

\begin{Entry}{里斯本大学}{7,12,5,3,8}{⾥,⽄,⽊,⼤,⼦}
  \begin{Phonetics}{里斯本大学}{li3si1ben3 da4xue2}
    \definition*{s.}{Universidade de Lisboa}
  \end{Phonetics}
\end{Entry}

\begin{Entry}{里程碑}{7,12,13}{⾥,⽲,⽯}
  \begin{Phonetics}{里程碑}{li3cheng2bei1}[][HSK 7-9]
    \definition[个,块,座]{s.}{marco; placas colocadas ao longo das estradas para indicar distâncias | ponto significativo no desenvolvimento; um evento muito importante na história; uma metáfora para um evento importante que pode servir como um marco no curso do desenvolvimento histórico}
  \end{Phonetics}
\end{Entry}

%%%%%%%%%% 重 %%%%%%%%%%
\subsection*{重}\addcontentsline{loh}{figure}{重}

\begin{Entry}{重}{9}{⾥}
  \begin{Phonetics}{重}{chong2}[][HSK 3]
    \definition*{s.}{Sobrenome: Chong}
    \definition{adv.}{novamente; mais uma vez}
    \definition{clas.}{usado para camadas}
    \definition{v.}{repetir; duplicar}
  \synonymref{沉}{chen2}
  \synonymref{浓}{nong2}
  \antonymref{轻}{qing1}
  \end{Phonetics}
  \begin{Phonetics}{重}{zhong4}[][HSK 1]
    \definition{adj.}{pesado; densidade elevada | profundo; sério; grau profundo | importante; significativo | discreto; prudente | considerável em quantidade ou valor}
    \definition[斤,公,斤,吨]{s.}{peso}
    \definition{v.}{colocar (colocar, pôr) ênfase em; dar valor a; atribuir importância a}
  \seealsoref{沉重}{chen2zhong4}
  \synonymref{沉}{chen2}
  \synonymref{沉重}{chen2zhong4}
  \synonymref{重新}{chong2xin1}
  \synonymref{浓}{nong2}
  \synonymref{深}{shen1}
  \antonymref{浅}{qian3}
  \antonymref{轻}{qing1}
  \end{Phonetics}
\end{Entry}

\begin{Entry}{重大}{9,3}{⾥,⼤}
  \begin{Phonetics}{重大}{zhong4da4}[][HSK 3]
    \definition{adj.}{excelente; importante; significativo; de grande importância}
  \end{Phonetics}
\end{Entry}

\begin{Entry}{重中之重}{9,4,3,9}{⾥,⼁,⼂,⾥}
  \begin{Phonetics}{重中之重}{zhong4zhong1zhi1zhong4}[][HSK 7-9]
    \definition{expr.}{o mais importante; de prioridade máxima; da maior importância}
  \synonymref{至关重要}{zhi4guan1-zhong4yao4}
  \antonymref{微不足道}{wei1bu4zu2dao4}
  \antonymref{无关紧要}{wu2guan1-jin3yao4}
  \antonymref{无足轻重}{wu2zu2-qing1zhong4}
  \end{Phonetics}
\end{Entry}

\begin{Entry}{重心}{9,4}{⾥,⼼}
  \begin{Phonetics}{重心}{zhong4xin1}[][HSK 7-9]
    \definition{s.}{Física: centro de gravidade | coração; núcleo; foco | ponto"-chave; o centro ou a parte principal da questão | baricentro; ponto mediano}
  \synonymref{核心}{he2xin1}
  \synonymref{焦点}{jiao1dian3}
  \synonymref{要点}{yao4dian3}
  \synonymref{中心}{zhong1xin1}
  \synonymref{中央}{zhong1yang1}
  \synonymref{重点}{zhong4dian3}
  \synonymref{主题}{zhu3ti2}
  \end{Phonetics}
\end{Entry}

\begin{Entry}{重申}{9,5}{⾥,⽥}
  \begin{Phonetics}{重申}{chong2shen1}[][HSK 7-9]
    \definition{v.}{reafirmar; reiterar}[他重申了自己对古巴的承诺。===Ele reiterou seu compromisso com Cuba.]
  \end{Phonetics}
\end{Entry}

\begin{Entry}{重任}{9,6}{⾥,⼈}
  \begin{Phonetics}{重任}{zhong4ren4}[][HSK 7-9]
    \definition{s.}{responsabilidade importante; grande responsabilidade; responsabilidade significativa; tarefa importante}
  \end{Phonetics}
\end{Entry}

\begin{Entry}{重伤}{9,6}{⾥,⼈}
  \begin{Phonetics}{重伤}{zhong4shang1}[][HSK 7-9]
    \definition{s.}{lesão grave | gravemente ferido}
  \end{Phonetics}
\end{Entry}

\begin{Entry}{重创}{9,6}{⾥,⼑}
  \begin{Phonetics}{重创}{zhong4chuang1}[][HSK 7-9]
    \definition{v.}{infligir grandes perdas | causar danos graves}
  \synonymref{严重}{yan2zhong4}
  \end{Phonetics}
\end{Entry}

\begin{Entry}{重合}{9,6}{⾥,⼝}
  \begin{Phonetics}{重合}{chong2he2}[][HSK 7-9]
    \definition{s.}{coincidir; sobrepor"-se}[两个假期的时间重合了。===Os dois feriados se sobrepõem.]
  \end{Phonetics}
\end{Entry}

\begin{Entry}{重阳节}{9,6,5}{⾥,⾩,⾋}
  \begin{Phonetics}{重阳节}{chong2yang2jie2}
    \definition*{s.}{Festa do Duplo Nove, Festival Yang, dia de subir aos lugares mais altos para evitar calamidades e dia do culto aos antepassados (9º dia do nono mês lunar)}
  \end{Phonetics}
\end{Entry}

\begin{Entry}{重返}{9,7}{⾥,⾡}
  \begin{Phonetics}{重返}{chong2fan3}[][HSK 7-9]
    \definition{v.}{retornar; voltar; de volta ao lugar original}
  \end{Phonetics}
\end{Entry}

\begin{Entry}{重建}{9,8}{⾥,⼵}
  \begin{Phonetics}{重建}{chong2jian4}[][HSK 6]
    \definition{s.}{restabelecimento; reconstrução}
    \definition{v.}{reconstruir; reconstruir; restabelecer; reabilitar}
  \end{Phonetics}
\end{Entry}

\begin{Entry}{重现}{9,8}{⾥,⾒}
  \begin{Phonetics}{重现}{chong2xian4}[][HSK 7-9]
    \definition{v.}{reaparecer | reproduzir}
  \end{Phonetics}
\end{Entry}

\begin{Entry}{重组}{9,8}{⾥,⽷}
  \begin{Phonetics}{重组}{chong2zu3}[][HSK 6]
    \definition{v.}{reestruturar; reorganizar; remanejar}
  \end{Phonetics}
\end{Entry}

\begin{Entry}{重视}{9,8}{⾥,⾒}
  \begin{Phonetics}{重视}{zhong4shi4}[][HSK 2]
    \definition{v.}{valorizar; dar peso a; atribuir importância a; prestar atenção a; considerar a virtude ou o talento de uma pessoa ou o papel de algo como importante e levá"-lo a sério}
  \synonymref{侧重}{ce4zhong4}
  \synonymref{崇尚}{chong2shang4}
  \synonymref{关心}{guan1xin1}
  \synonymref{讲究}{jiang3jiu5}
  \synonymref{看重}{kan4zhong4}
  \synonymref{珍惜}{zhen1xi1}
  \synonymref{珍重}{zhen1zhong4}
  \synonymref{注意}{zhu4/yi4}
  \synonymref{注重}{zhu4zhong4}
  \antonymref{鄙视}{bi3shi4}
  \antonymref{不屑}{bu2xie4}
  \antonymref{忽略}{hu1lve4}
  \antonymref{忽视}{hu1shi4}
  \antonymref{回避}{hui2bi4}
  \antonymref{歧视}{qi2shi4}
  \antonymref{无视}{wu2shi4}
  \antonymref{小看}{xiao3kan4}
  \end{Phonetics}
\end{Entry}

\begin{Entry}{重型}{9,9}{⾥,⼟}
  \begin{Phonetics}{重型}{zhong4xing2}[][HSK 7-9]
    \definition{adj.}{para serviço pesado; pesado | serviço pesado | grande calibre}
  \antonymref{轻型}{qing1xing2}
  \end{Phonetics}
\end{Entry}

\begin{Entry}{重复}{9,9}{⾥,⼢}
  \begin{Phonetics}{重复}{chong2fu4}[][HSK 2]
    \definition{v.}{repetir; iterar; duplicar; reduplicar | fazer algo novamente; repetir as mesmas palavras, fazer as mesmas coisas}
  \seealsoref{反复}{fan3fu4}
  \synonymref{反复}{fan3fu4}
  \synonymref{屡次}{lv3ci4}
  \synonymref{循环}{xun2huan2}
  \synonymref{再三}{zai4san1}
  \antonymref{不再}{bu2zai4}
  \antonymref{唯一}{wei2yi1}
  \end{Phonetics}
\end{Entry}

\begin{Entry}{重点}{9,9}{⾥,⽕}
  \begin{Phonetics}{重点}{chong2dian3}
    \definition[个]{adj./adv./s.}{nota principal; ponto"-chave; ponto focal; ênfase}
  \end{Phonetics}
  \begin{Phonetics}{重点}{zhong4dian3}[][HSK 2]
    \definition[个]{s.}{nota principal; ponto"-chave; ponto}
  \synonymref{关键}{guan1jian4}
  \synonymref{核心}{he2xin1}
  \synonymref{要点}{yao4dian3}
  \synonymref{中心}{zhong1xin1}
  \synonymref{重心}{zhong4xin1}
  \synonymref{主要}{zhu3yao4}
  \synonymref{着重}{zhuo2zhong4}
  \antonymref{无足轻重}{wu2zu2-qing1zhong4}
  \antonymref{一般}{yi4ban1}
  \end{Phonetics}
\end{Entry}

\begin{Entry}{重要}{9,9}{⾥,⾑}
  \begin{Phonetics}{重要}{zhong4yao4}[][HSK 1]
    \definition{adj.}{importante; significativo; relevante; de grande importância, função e impacto}
  \synonymref{关键}{guan1jian4}
  \synonymref{首要}{shou3yao4}
  \synonymref{要紧}{yao4jin3}
  \synonymref{重大}{zhong4da4}
  \synonymref{主要}{zhu3yao4}
  \antonymref{普通}{pu3tong1}
  \antonymref{一般}{yi4ban1}
  \end{Phonetics}
\end{Entry}

\begin{Entry}{重重}{9,9}{⾥,⾥}
  \begin{Phonetics}{重重}{chong2chong2}
    \definition{adj.}{camada sobre camada; descrever coisas como numerosas e densas}
  \synonymref{层层}{ceng2ceng2}
  \antonymref{单一}{dan1yi1}
  \end{Phonetics}
  \begin{Phonetics}{重重}{zhong4zhong4}
    \definition{adv.}{pesadamente; severamente}
  \end{Phonetics}
\end{Entry}

\begin{Entry}{重逢}{9,10}{⾥,⾡}
  \begin{Phonetics}{重逢}{chong2feng2}
    \definition{v.}{reunir; encontrar"-se novamente; reencontrar"-se após uma longa separação; rever"-se}
  \end{Phonetics}
\end{Entry}

\begin{Entry}{重量}{9,12}{⾥,⾥}
  \begin{Phonetics}{重量}{zhong4liang4}[][HSK 4]
    \definition[个]{s.}{peso; a magnitude da força da gravidade em um objeto}
  \end{Phonetics}
\end{Entry}

\begin{Entry}{重量级}{9,12,6}{⾥,⾥,⽷}
  \begin{Phonetics}{重量级}{zhong4liang4ji2}[][HSK 7-9]
    \definition{s.}{peso pesado (boxe, etc.)}
  \end{Phonetics}
\end{Entry}

\begin{Entry}{重叠}{9,13}{⾥,⼜}
  \begin{Phonetics}{重叠}{chong2die2}[][HSK 7-9]
    \definition{s.}{reduplicação; gramaticalmente, refere"-se ao uso da mesma palavra ou frase duas vezes}
    \definition{v.}{sobrepor; por um em cima do outro; colocar as mesmas coisas juntas camada por camada}
  \end{Phonetics}
\end{Entry}

\begin{Entry}{重新}{9,13}{⾥,⽄}
  \begin{Phonetics}{重新}{chong2xin1}[][HSK 2]
    \definition{adv.}{novamente; de novo; significa repetir uma ação ou comportamento já realizado | indica que se deve começar do início (mudança de método ou conteúdo)}
  \seealsoref{重}{chong2}
  \seealsoref{再}{zai4}
  \synonymref{重}{chong2}
  \synonymref{重复}{chong2fu4}
  \synonymref{从头}{cong2tou2}
  \synonymref{再}{zai4}
  \synonymref{再次}{zai4ci4}
  \synonymref{再度}{zai4du4}
  \antonymref{保持}{bao3chi2}
  \antonymref{不再}{bu2zai4}
  \antonymref{继续}{ji4xu4}
  \antonymref{维持}{wei2chi2}
  \antonymref{依旧}{yi1jiu4}
  \end{Phonetics}
\end{Entry}

\begin{Entry}{重播}{9,15}{⾥,⼿}
  \begin{Phonetics}{重播}{chong2bo1}[][HSK 7-9]
    \definition{v.}{retransmitir (um programa de rádio ou TV) | Agricultura: ressemear (o mesmo campo)}
  \end{Phonetics}
\end{Entry}

%%%%%%%%%% 野 %%%%%%%%%%
\subsection*{野}\addcontentsline{loh}{figure}{野}

\begin{Entry}{野}{11}{⾥}
  \begin{Phonetics}{野}{ye3}[][HSK 6]
    \definition*{s.}{Sobrenome: Ye}
    \definition{adj.}{(de plantas ou animais) selvagem; incultivado; não domesticado; indomável (opp. 家) | rude; áspero | desenfreado; abandonado; indisciplinado | ilícito; sem licença}
    \definition{s.}{espaço aberto; o aberto | limite; fronteira | não está no poder; fora do cargo}
  \seealsoref{家}{jia1}
  \end{Phonetics}
\end{Entry}

\begin{Entry}{野心}{11,4}{⾥,⼼}
  \begin{Phonetics}{野心}{ye3xin1}[][HSK 7-9]
    \definition{s.}{carreirismo; ambição desmedida; plano insensato; um desejo enorme e excessivo por território, poder ou fama}
  \synonymref{打算}{da3suan5}
  \synonymref{计划}{ji4hua4}
  \synonymref{盘算}{pan2suan5}
  \synonymref{企图}{qi3tu2}
  \synonymref{妄想}{wang4xiang3}
  \synonymref{希望}{xi1wang4}
  \synonymref{阴谋}{yin1mou2}
  \end{Phonetics}
\end{Entry}

\begin{Entry}{野外}{11,5}{⾥,⼣}
  \begin{Phonetics}{野外}{ye3wai4}[][HSK 7-9]
    \definition{s.}{campo; ar livre; campo aberto; lugares distantes de áreas residenciais}
  \synonymref{郊外}{jiao1wai4}
  \synonymref{旷野}{kuang4ye3}
  \end{Phonetics}
\end{Entry}

\begin{Entry}{野生}{11,5}{⾥,⽣}
  \begin{Phonetics}{野生}{ye3sheng1}[][HSK 6]
    \definition{adj.}{selvagem; não cultivado; não domesticado}
  \end{Phonetics}
\end{Entry}

\begin{Entry}{野炊}{11,8}{⾥,⽕}
  \begin{Phonetics}{野炊}{ye3chui1}[][HSK 7-9]
    \definition{v.}{cozinhar ao ar livre; cozinhar sobre uma fogueira na natureza selvagem}
  \end{Phonetics}
\end{Entry}

\begin{Entry}{野兽}{11,11}{⾥,⼋}
  \begin{Phonetics}{野兽}{ye3shou4}[][HSK 7-9]
    \definition[头]{s.}{besta; animal selvagem; mamíferos que não sejam animais de criação}
  \antonymref{美女}{mei3nv3}
  \end{Phonetics}
\end{Entry}

\begin{Entry}{野营}{11,11}{⾥,⾋}
  \begin{Phonetics}{野营}{ye3ying2}[][HSK 7-9]
    \definition{v.}{acampar; ficar em um acampamento temporário sem abrigo; acampar e praticar certas atividades na natureza}
  \end{Phonetics}
\end{Entry}

\begin{Entry}{野蛮}{11,12}{⾥,⾍}
  \begin{Phonetics}{野蛮}{ye3man2}[][HSK 7-9]
    \definition{adj.}{selvagem; incivilizado; inculto | cruel; brutal; bárbaro; desumano}
  \synonymref{残忍}{can2ren3}
  \synonymref{粗暴}{cu1bao4}
  \synonymref{粗鲁}{cu1lu3}
  \synonymref{凶恶}{xiong1'e4}
  \antonymref{仁慈}{ren2ci2}
  \antonymref{温柔}{wen1rou2}
  \antonymref{文明}{wen2ming2}
  \end{Phonetics}
\end{Entry}

\begin{Entry}{野餐}{11,16}{⾥,⾷}
  \begin{Phonetics}{野餐}{ye3can1}[][HSK 7-9]
    \definition[个,次]{s.}{piquenique; diversas opções de comida para refeições ao ar livre}
    \definition{v.}{fazer um piquenique; ir a um piquenique; comer em meio à natureza selvagem}
  \end{Phonetics}
\end{Entry}

%%%%%%%%%% 量 %%%%%%%%%%
\subsection*{量}\addcontentsline{loh}{figure}{量}

\begin{Entry}{量}{12}{⾥}
  \begin{Phonetics}{量}{liang2}[][HSK 4]
    \definition{v.}{medir | estimar; dimensionar}
  \end{Phonetics}
  \begin{Phonetics}{量}{liang4}
    \definition{s.}{instrumento de medida; antigamente, o termo se referia a objetos como baldes e litros, que medem o volume | capacidade (para tolerância ou ingestão de alimentos ou bebidas); refere"-se ao limite do que pode ser acomodado | quantidade; valor; volume; número}
    \definition{v.}{estimar; medir; pesar}
  \end{Phonetics}
\end{Entry}

%%%%% EOF %%%%%

